\newpage
\section{Results and Analysis}\label{sec:result}

\subsection{Comparison of the Four Classifiers}
All four classifiers were trained on the same feature set and evaluated with LOSO cross-validation over the 3 group members. Table~\ref{tab:model_comparison} summarizes the overall mean accuracy and its standard deviation across the three folds.

\begin{table}[!htb]
    \centering
    \caption{LOSO cross-validation accuracy of the four classifiers}
    \label{tab:model_comparison}
    \begin{tabularx}{\textwidth}{@{} l l l X @{}}
        \toprule
        \textbf{Model} & \textbf{Mean Accuracy} & \textbf{Std.} & \textbf{Remarks} \\
        \midrule
        MLP & \todo{} & \todo{$\pm$ } & \todo{} \\
        SVM & \todo{} & \todo{$\pm$ } & \todo{} \\
        RF  & \todo{} & \todo{$\pm$ } & \todo{} \\
        \todo{Self-selected} & \todo{} & \todo{$\pm$ } & \todo{} \\
        \bottomrule
    \end{tabularx}
\end{table}

\noindent \todo{State which model performed best and by how much.}

\subsection{Effect of Preprocessing and Feature Engineering}
To quantify how much of the performance comes from the processing pipeline rather than from the classifier itself, we compare a reference configuration against our final pipeline. \todo{Define your own reference configuration, e.g., the raw time-domain samples fed directly to the classifiers with no filtering and report its accuracy. If you performed a step-by-step ablation, present it here.}

\begin{table}[!htb]
    \centering
    \caption{Ablation study of the processing pipeline (accuracy of the best model)}
    \label{tab:ablation}
    \begin{tabularx}{\textwidth}{@{} l l X @{}}
        \toprule
        \textbf{Configuration} & \textbf{Mean Accuracy} & \textbf{Observation} \\
        \midrule
        \todo{Raw time-domain samples} & \todo{} & \todo{} \\
        \todo{+ Bandpass filtering} & \todo{} & \todo{} \\
        \todo{+ ...} & \todo{} & \todo{} \\
        \todo{Final pipeline} & \todo{} & \todo{} \\
        \bottomrule
    \end{tabularx}
\end{table}

\subsection{Detailed Results of the Best Model}
\todo{Insert your final results image here.}
\begin{figure}[h!]
    \centering
    % \includegraphics[width=\linewidth]{your_final_results.png}
    \caption{LOSO cross-validation results of the best model, including the accuracy of the four models, the per-subject accuracy for the 3 group members, and the overall $3 \times 3$ confusion matrix.}
    \label{fig:final_results}
\end{figure}

\noindent \todo{Insert your final terminal output summary here.}
\begin{verbatim}
Summary of Classification Results:
 - MLP         : 0.XXX +/- 0.XXX
 - SVM         : 0.XXX +/- 0.XXX
 - RF          : 0.XXX +/- 0.XXX
 - SelfSelected: 0.XXX +/- 0.XXX

Best Performing Model: XXX with Accuracy 0.XXX

Detailed Class Report for Best Model [XXX]:
[Relax Class]
  - Accuracy (Recall): 0.XXX
  - Precision: 0.XXX

[Focus Class]
  - Accuracy (Recall): 0.XXX
  - Precision: 0.XXX

[Stroop Class]
  - Accuracy (Recall): 0.XXX
  - Precision: 0.XXX
\end{verbatim}

\subsection{Result Analysis and Discussion}
\todo{Analyze your results in depth.}
\begin{itemize}[leftmargin=*]
    \item \textbf{Accuracy Comparison:} Compare the four models in Table~\ref{tab:model_comparison}. Did the best model reach the $\geq 75\%$ target? Explain why the winning model suits this dataset.
    \item \textbf{Confusion Matrix Analysis:} Using the $3 \times 3$ confusion matrix in Figure~\ref{fig:final_results}:
    \begin{itemize}
        \item Which class was the easiest for the model to identify? Relate your answer to the expected spectral signature of that state at Fp1.
        \item Which two classes were most frequently confused with each other?
        \item How did your feature engineering or post-processing help resolve this confusion?
    \end{itemize}
    \item \textbf{Inter-Subject Variability:} Looking at the per-subject bar chart in Figure~\ref{fig:final_results}, did the model perform equally well on all 3 group members? If one member had a significantly lower accuracy, discuss potential reasons.
    \item \textbf{Method Effectiveness Analysis:} Based on Table~\ref{tab:ablation}, discuss which processing step contributed the most to the improvement. Justify your reasoning.
    \item \textbf{Limitations:} \todo{Discuss the limitations of a single-channel Fp1 recording with only 3 subjects, and what would be needed for a more reliable subject-independent conclusion.}
\end{itemize}